% ============================================================
% REVISED INTRODUCTION
% Structure:
%   1. Problem motivation
%   2-5. Prior work with logical transitions (each paragraph leads to next)
%   6. "In general, several critical gaps remain..."
%   7. "In this research, we propose..."
%   8. Experimental setup + key results summary
%   9. Main contributions list
%  10. Paper organization
% ============================================================
\section{Introduction}
\label{sec:intro}

\IEEEPARstart{T}{he} explosive growth of IoT deployments---spanning smart
cities, industrial sensing, and indoor positioning---has introduced
unprecedented demands on wireless network infrastructure. Devices operating
across heterogeneous technologies such as Sigfox, LoRaWAN, and Indoor WiFi
continuously generate high-dimensional telemetry streams that must be
processed, monitored, and acted upon with minimal latency and energy
overhead~\cite{ref_iot_survey}. Two fundamental challenges arise: the raw
feature space is high-dimensional and noisy, making downstream learning
inefficient; and the network state evolves continuously over time, requiring
adaptive mechanisms that simultaneously classify connection quality and detect
anomalous behaviour.

A natural approach to the dimensionality challenge is to extract compact,
task-relevant representations from raw sensor observations before any
downstream processing. Liu et al.~\cite{ref_vae_wireless} proposed a
task-adaptive semantic compression scheme employing a Variational Autoencoder
(VAE) to encode raw sensor observations into compact representations,
substantially reducing communication overhead in resource-constrained
networks. The compression rate adapts to instantaneous channel quality,
yielding strong reconstruction fidelity across diverse wireless conditions.
However, this scheme targets static point-to-point links without temporal
state tracking, provides no anomaly-detection capability, and does not
coordinate multiple agents---restricting its applicability to the dynamic,
heterogeneous IoT settings addressed in this work.

Beyond feature compression, the problem of adaptive resource management under
time-varying interference has been studied through deep reinforcement learning
(DRL). Hazarika et al.~\cite{ref_marl_wireless} developed a deep RL framework
for computation offloading in Internet-of-Vehicles (IoV) networks, training
DQN agents on raw network-state features to achieve adaptive resource
management with low control overhead. A key limitation, however, is that
raw, high-dimensional feature vectors serve as the agent state without
semantic compression, inflating the effective state space and degrading
sample efficiency. Moreover, no channel-quality feedback enters the reward
signal, preventing the agent from distinguishing nominal from degraded link
conditions. These shortcomings motivate coupling feature compression with a
predictive channel model before resource allocation.

To address the channel-awareness limitation, digital twin (DT) technology has
been introduced as a mechanism for real-time channel prediction and quality
monitoring. Hazarika et al.~\cite{ref_dt_uav} integrated a digital twin with
a hybrid machine learning classifier for resource allocation in UAV-aided IoV
networks, enabling one-step-ahead channel prediction that yields measurable
gains over purely reactive baselines. Nevertheless, quality classification
relies on fixed decision thresholds that cannot adapt to heterogeneous dataset
statistics, anomaly detection is absent, and raw features are fed directly to
the allocation module without semantic compression---leaving the system unable
to flag silent sensor faults or generalize across heterogeneous channel
conditions.

Recognizing the need for a more tightly integrated system, Hazarika
et al.~\cite{ref_semadt} combined semantic compression, predictive digital
twin forecasting, and multi-agent deep RL with a graph neural network (GNN)
for maritime search-and-rescue IoT operations, demonstrating that semantic
state inputs and DT-predicted quality signals together substantially outperform
non-predictive MARL baselines. Despite this advance, the GCN layer applies
symmetric neighborhood aggregation without capturing heterogeneous inter-agent
relationship strengths, whereas graph attention networks~\cite{ref_gat_comm}
do so via learned attention weights. Furthermore, the anomaly detector relies
on a fixed scalar threshold rather than a statistically principled multivariate
test~\cite{ref_anomaly_survey,ref_lof}, and the evaluation is confined to a
single maritime domain, leaving cross-technology generalizability unexplored.

In general, several critical gaps remain in the existing literature. First,
semantic compression, channel state monitoring, and resource allocation are
treated as independent modules with no unified feedback path between them.
Second, existing digital twin implementations address either quality
classification \emph{or} anomaly detection in isolation, rather than jointly
performing both within a single principled state estimator. Third, no prior
system has been validated across heterogeneous IoT technologies in a unified
architecture, limiting confidence in cross-domain generalizability.

In this research, we propose a unified three-stage intelligent framework to
simultaneously address the above gaps. The framework integrates: (i) a
Variational Autoencoder (VAE) for semantic compression of raw IoT features
into compact latent representations; (ii) an Extended Kalman Filter
(EKF)-based Digital Twin for joint connection quality classification and
anomaly detection via the Normalized Innovation Squared (NIS) test; and
(iii) a Multi-Agent Deep Reinforcement Learning module augmented with a
Graph Attention Network (MADRL-GAT) for cooperative resource allocation
guided by the semantic state and quality feedback produced in the earlier
stages.

We evaluate the proposed framework on three real-world IoT datasets spanning
heterogeneous wireless technologies: Sigfox Antwerp, LoRaWAN Italy, and
Indoor WiFi. Experimental results show that the EKF Digital Twin achieves
quality classification AUC of up to 0.972, outperforming K-Means by up to
80.4\% on the LoRaWAN dataset, while attaining anomaly detection AUC of
0.987 on the Indoor dataset. The MADRL-GAT module converges stably within
50 training episodes, yielding evaluation rewards of 1,916, 2,519, and 3,089
across the three datasets, respectively, confirming that the integrated
semantic and quality-aware state representation accelerates policy learning.

The main contributions of this paper are summarized as follows:
\begin{itemize}
\item \textbf{We propose a unified three-stage framework} that tightly couples
VAE semantic compression, EKF Digital Twin monitoring, and MADRL-GAT resource
allocation into a single pipeline, where compressed latent codes serve as both
the MADRL state space and the input for quality-aware reward computation.
\item \textbf{We design an EKF-based Digital Twin that simultaneously performs
connection quality classification and anomaly detection} within a single filter
pass, using a data-adaptive Youden's J threshold on the estimated RSSI for
quality labeling and the NIS chi-squared test for fault flagging.
\item \textbf{We validate the framework across three heterogeneous real-world
IoT datasets}---Sigfox Antwerp, LoRaWAN Italy, and Indoor WiFi---demonstrating
consistent robustness across technologies with substantially different channel
characteristics and feature dimensionalities.
\end{itemize}

The remainder of this paper is organised as follows. Section~\ref{sec:model}
presents the system model. Section~\ref{sec:problem} formulates the
optimization problem. Section~\ref{sec:method} details the proposed framework.
Section~\ref{sec:experiments} reports experimental results.
Section~\ref{sec:conclusion} concludes the paper.
